\section{Transporte hasta el coeficiente crítico de la máscara}
Fijemos los $v$ sitios $U$ antes de escoger las activaciones y el exterior común
obtenido al quitarlos de la máscara de torres. Sea $c$ el número de componentes
de la incidencia completa y supongamos que el grafo tiene exactamente un ciclo.
Usamos la descomposición del exterior de \emph{Patrones finitos}, parte 03, en la
misma rama para todas las caras. Al coste cero quedan una constante común y campos
afines, absorbidos en las actividades de sitio. El término de reacción $R''$
tiene coste positivo después de esa extracción. Esto se ve expandiendo
$s^{-1}R_0(sy)$: sus términos de grados cero y uno en $y$ son constante y afín.
Los términos de coste cero de las integrales $E_{ab}$ y del prefactor
son constantes comunes, y el resto puro de torres comienza en $s$.

\paragraph{Cuáles son las derivadas relevantes.}
Póngase $s=1/j$ en la evaluación física. A coste $w=c-1$, un término conexo
debe alcanzar los $c$ grupos para sobrevivir a la diferencia de activaciones.
La cuenta de incidencias $d=1+w-\sum_a(l_a+1-h_a)-L_I$ obliga a déficit cero
y $L_I=0$. Aquí $h_a$ cuenta los grupos del término $a$ de coste $l_a$.
Sólo quedan árboles de interacciones de acción, multilineales entre grupos distintos.
Si $c=1$, el coste es cero y se aplica directamente la identidad local.

Se enraíza el árbol en el grupo con ciclo. En una interacción hoja de amplitud $t$
escribamos $V=L_0\prod_{a\ge1}L_a$, con el grupo de enlace en el núcleo y una
forma lineal por cada grupo privado. Al anular $t$ estos últimos son independientes
del núcleo, aun con las demás amplitudes del árbol encendidas. Sus medias exactas
$m_a$ coinciden con las de la silla, pues las hojas son forestales.
Para la diferencia de logaritmos $\Delta$ se tiene
\[
 \left.\partial_t\Delta\right|_{t=0}
    =\left(\prod_{a\ge1}m_a\right)D_0\Delta_{\rm nucleo},
\]
donde $D_0$ deriva en el campo que acopla $L_0$. Para el exacto se usa la
factorización. Para la silla se deriva el valor estacionario. En el cofactor,
la Hessiana explícita $V''$ sólo tiene bloques cruzados y su traza contra la inversa
por bloques se anula. El desplazamiento de la silla de cada rama equivale a un
cambio de su campo; su prefactor completo vale uno en todos los campos de la rama.
La variación en el núcleo es la de $D_0$ por el producto de las medias privadas.
Se mantienen la variación de sus factores estáticos y la de su determinante.
Esta derivada en amplitud cero extrae el término lineal de la interacción marcada.
La repetición de una interacción en un árbol crítico se trata con la cuenta de
incidencias: dos copias que comparten al menos dos grupos formarían un ciclo.
Al podar sucesivamente se aplica también la regla del producto a las medias anteriores.

\paragraph{Divisibilidad después de las caras.}
En una cara que conserva el ciclo sus sitios son siempre los mismos.
En la rama local $x_i=\zeta_i b_i(\zeta)$ con $b_i$ holomorfo, $1-x_i$ invertible.
La diferencia local de la parte 01 pertenece al ideal
$I_{\mathcal C}=(\prod_{i\in\mathcal C}\zeta_i^2)$. Cada derivada de campo es
una combinación de $\zeta_i\partial_{\zeta_i}$, que preserva ese ideal.
También lo preservan los factores holomorfos de las ramas. Las caras que rompen
el ciclo tienen raíz forestal con diferencia crítica cero una vez eliminados
los términos que no alcanzan todos los grupos.

Sobre las caras, las constantes del exterior y el modo nulo RC son los mismos.
Al apagar un sitio, la columna correspondiente de $(-H_{\rm B}''-R'')V_0$
es en el límite el vector coordenado; los factores escalares se restringen a
la misma cara. Esto es la identidad de cofactor de la parte 07 citada.
La partición exacta o su aproximante por átomos también pierde ese sitio.
La diferencia completa de activaciones equivale entonces a su proyección
$\prod_{i\in U}(I-\operatorname{ev}_{\zeta_i=0})$. Esta operación
y la extracción de $[s^w]$ conmutan para los aproximantes holomorfos comunes.
Así su coeficiente proyectado de coste $w=c-1$ es divisible por
\[
                 \prod_{i\in U}\zeta_i\ \prod_{i\in\mathcal C}\zeta_i .
\]
La afirmación de coeficientes se hace sobre el aproximante. El error físico real
en el tamaño entero se trata con su cota propia, sin diferenciarlo en $s$.
